\documentclass[14pt,a4paper]{extarticle}
\usepackage{fontspec}
\usepackage{amsmath}
\usepackage{geometry}
\geometry{left=30mm,right=10mm,top=20mm,bottom=20mm}
\usepackage{setspace}
\usepackage{array}
\usepackage{longtable}
\usepackage{booktabs}
\usepackage{ragged2e}
\usepackage{enumitem}
\usepackage{xcolor}

\setmainfont{Times New Roman}
\onehalfspacing
\setlength{\parindent}{1.25em}
\sloppy
\tolerance=2000
\emergencystretch=3em
\pagestyle{plain}

\newcommand{\code}[1]{\texttt{#1}}
\newcommand{\tabcap}[1]{\par\smallskip\noindent\textbf{#1}\par\smallskip}

\begin{document}

%======================= TITLE PAGE =======================
\thispagestyle{empty}
\begin{center}
\vspace*{2.5cm}
{\large Проект \code{physics\_md\_al\_fe}}\\[0.4cm]
{\large Научно-технический отчёт о результатах численного моделирования}\\[2.2cm]

{\LARGE\bfseries Отчёт о численном моделировании деформации алюминиевой матрицы с интерметаллидным включением при магнитострикционном воздействии}\\[2.2cm]

{\large Этап Stage B (focused, 100k атомов) — завершён;\\[0.2cm]
этап Stage C (1M атомов) — запущен}\\[1.6cm]

{\large Метод: молекулярная динамика (LAMMPS, потенциал MEAM, GPU KOKKOS/CUDA)}\\[2.6cm]

\vfill
{\large 17 июня 2026 г.}
\end{center}
\newpage

%======================= ABSTRACT =======================
\section*{Аннотация}
В работе методом классической молекулярной динамики исследовано влияние магнитострикционной деформации интерметаллидного включения Fe$_4$Al$_{13}$ на дефектную структуру окружающей алюминиевой матрицы. Магнитное поле не моделируется явно: его действие задаётся как объёмосохраняющая собственная деформация (eigenstrain) атомов включения. Рассмотрена реалистичная конфигурация — включение у границы зерна алюминия с вакансиями в матрице. Выполнены два расчётных варианта по 100\,000 шагов: физический ($\varepsilon_z = 0.0025$) и усиленный ($\varepsilon_z = 0.0100$). В обоих вариантах анализ дислокаций (DXA) не выявил подтверждённых дислокационных линий; доля атомов в HCP-окружении пренебрежимо мала и не растёт с деформацией, а высокая доля атомов типа OTHER объясняется структурным фоном границы зерна, интерфейса включения и тепловыми смещениями при 300\,K. Сделан вывод о недостаточности масштаба 100k для зарождения дислокаций и обоснован переход к масштабному расчёту $\approx$\,1 млн атомов (Stage C), который на момент составления отчёта запущен.

%======================= TOC =======================
\section*{Содержание}
\begin{enumerate}[leftmargin=2.2em,itemsep=0pt]
\item Введение
\item Постановка задачи
\item Исходные данные и допущения
\item Методика расчёта
\item Расчётная модель
\item Численная реализация
\item Результаты расчётов Stage B 100k
\item Анализ результатов
\item План Stage C 1M
\item Ограничения
\item Заключение
\item Список источников
\item Приложения
\end{enumerate}
\newpage

%======================= 1 =======================
\section{Введение}
В алюминиевых сплавах технического назначения часто присутствуют железосодержащие интерметаллидные фазы, в частности Fe$_4$Al$_{13}$ (Al$_{13}$Fe$_4$). При магнитной обработке материала такие включения могут испытывать магнитострикционную деформацию и передавать в матрицу локальные напряжения. Эти напряжения способны вызывать локальную пластическую перестройку и зарождение дислокаций вблизи включения. Настоящий отчёт описывает результаты атомистического моделирования этого механизма на наномасштабе. Основной вопрос: при каком масштабе модели и уровне воздействия в алюминиевой матрице у включения возникают \textbf{подтверждённые} дефекты и дислокации.

%======================= 2 =======================
\section{Постановка задачи}
Требуется средствами молекулярной динамики: построить модель «алюминиевая матрица + эллипсоидное включение Fe$_4$Al$_{13}$» в реалистичной конфигурации (включение у границы зерна, вакансии в матрице); задать механический эквивалент магнитострикционного воздействия через собственную деформацию включения; провести расчёты при физическом и усиленном уровне воздействия; количественно оценить дефектную структуру матрицы (дислокации, доли FCC/HCP/OTHER, зону пластической перестройки); по результатам обосновать целесообразность масштабного расчёта.

Критерий подтверждённой дислокации: \code{dislocation\_segments} $> 0$, либо \code{dislocation\_line\_length\_A} $> 0$, либо \code{dislocation\_density\_per\_m2} $> 0$.

%======================= 3 =======================
\section{Исходные данные и допущения}
\textbf{Физические исходные данные:} матрица — алюминий, ГЦК-решётка, параметр $a = 4.05$\,\AA; включение — Fe$_4$Al$_{13}$ (Al$_{13}$Fe$_4$), моноклинная фаза (C2/m); температура расчёта $T = 300$\,K; межатомное взаимодействие — потенциал MEAM (трек \code{MEAM\_Jelinek\_2012}).

\textbf{Магнитное воздействие.} В проектной постановке рассматривается поле порядка $B \approx 0.7$\,Тл. Однако в расчётах поле \textbf{не задаётся явно}: управляющим параметром служит собственная деформация включения $\varepsilon_z$. Значения $\varepsilon_z = 0.0025$ и $\varepsilon_z = 0.0100$ выбраны как физически близкий и усиленный варианты.

\textbf{Допущения:} магнитный эффект вводится механически, без явного спинового моделирования; собственная деформация прикладывается только к атомам включения; деформация объёмосохраняющая; рассматривается наномасштабный фрагмент; per-atom напряжения трактуются как сравнительные вириальные величины.

%======================= 4 =======================
\section{Методика расчёта}
\subsection{Связь магнитострикции, напряжения и деформации}
Оценочная связь воздействия с напряжением и деформацией:
\begin{equation}
\sigma_m = \lambda_m \, B
\end{equation}
где $\lambda_m$ — эффективная магнитострикционная постоянная, $B$ — индукция поля.
\begin{equation}
\varepsilon \approx \sigma / E
\end{equation}
где $\sigma$ — действующее напряжение, $E$ — модуль Юнга материала включения. Из (1)–(2) видно, что характерная деформация включения при умеренных полях мала (доли процента), что и отражено выбором $\varepsilon_z = 0.0025$ и $\varepsilon_z = 0.0100$.

\subsection{Реализация собственной деформации (eigenstrain)}
Атомам включения задаётся объёмосохраняющая собственная деформация относительно центра включения:
\begin{equation}
\varepsilon_x = \varepsilon_y = -\varepsilon_z/2, \qquad \varepsilon_z > 0
\end{equation}
Так, для $\varepsilon_z = 0.0025$: $\varepsilon_x = \varepsilon_y = -0.00125$; для $\varepsilon_z = 0.0100$: $\varepsilon_x = \varepsilon_y = -0.005$. Деформация переводит включение в состояние с осевым растяжением и поперечным сжатием, создавая локальное напряжение на границе включение–матрица.

\subsection{Анализ дефектов}
Структурный анализ выполняется средствами OVITO: классификация локальной структуры (CNA/PTM) — доли FCC, HCP, OTHER; анализ дислокаций (DXA) — число сегментов, суммарная длина линий, плотность; индикатор дефектов упаковки — HCP-атомы внутри ГЦК-матрицы; протяжённость зоны пластичности — дефектные атомы за пределами нормированной оболочки 1.3 от интерфейса включения.

%======================= 5 =======================
\section{Расчётная модель}
\tabcap{Таблица 1. Параметры расчётной модели (focused Stage B 100k).}
\begin{longtable}{p{0.46\textwidth} p{0.46\textwidth}}
\toprule
Параметр & Значение \\
\midrule
\endhead
Размер ячейки & $105.3 \times 105.3 \times 157.95$\,\AA \\
Всего атомов & 103\,037 \\
Атомов матрицы (Al) & 98\,788 \\
Атомов включения (Fe$_4$Al$_{13}$) & 4\,249 \\
Полуоси включения & $19.5 \times 19.5 \times 39.0$\,\AA\ (1:1:2) \\
Положение & у границы зерна \\
Разориентация границы & $36.87^{\circ}$ (плоскость по оси $x$) \\
Зазор включение–граница & $\approx 7.2$\,\AA \\
Предефекты & 200 вакансий, оболочка 4\,\AA \\
Температура & 300\,K \\
Потенциал & MEAM (\code{MEAM\_Jelinek\_2012}) \\
\bottomrule
\end{longtable}

Два варианта собственной деформации включения: физический $\varepsilon_z = 0.0025$ ($\varepsilon_x = \varepsilon_y = -0.00125$) и усиленный $\varepsilon_z = 0.0100$ ($\varepsilon_x = \varepsilon_y = -0.005$).

%======================= 6 =======================
\section{Численная реализация}
Расчёт выполняется пакетом LAMMPS с ускорением на GPU (KOKKOS/CUDA, парный стиль \code{meam/kk}), на видеокарте NVIDIA RTX 3060 (sm\_86, CUDA 12.4).

\tabcap{Таблица 2. Цепочка расчёта и параметры стадий.}
\begin{longtable}{p{0.18\textwidth} r p{0.56\textwidth}}
\toprule
Стадия & Шагов & Назначение \\
\midrule
\endhead
prep & 8\,000 & прогрев 50$\to$300\,K (3\,000) + NVT 300\,K (5\,000), без минимизации \\
smoke & 2\,000 & проверка стабильности и раннего сигнала \\
production & 100\,000 & основной расчёт, порциями по 10\,000 шагов \\
\bottomrule
\end{longtable}

Особенности реализации: порционная production с записью \code{restart} и проверкой лога после каждой порции; resume с последнего корректного restart; watchdog/hang recovery (контроль зависания по CPU-времени и росту файлов, снятие зависшей порции и повтор один раз); обход проблемы соседних списков на GPU \code{neigh\_modify delay 0 every 10 check no}. Производительность: $\approx 2.69$ шага/с при 103\,037 атомах, $\approx 10.3$\,ч на вариант production.

%======================= 7 =======================
\section{Результаты расчётов Stage B 100k}
Оба варианта завершены штатно (100\,000/100\,000 шагов, код выхода 0, все порции с первой попытки, без зависаний).

\tabcap{Таблица 3. Результаты focused Stage B 100k.}
\begin{longtable}{p{0.50\textwidth} r r}
\toprule
Показатель & eps0025 & eps0100 \\
\midrule
\endhead
$\varepsilon_z$ & 0.0025 & 0.0100 \\
Число атомов & 103\,037 & 103\,037 \\
Число шагов & 100\,000 & 100\,000 \\
DXA segments & 0 & 0 \\
DXA length, \AA & 0.0 & 0.0 \\
DXA density, 1/м$^2$ & 0.0 & 0.0 \\
HCP atoms (матрица) & 8 & 1 \\
HCP beyond shell & 3 & 0 \\
OTHER atoms (матрица) & 17\,200 & 16\,144 \\
defect beyond shell & 13\,445 & 12\,517 \\
Вывод & дислокаций нет & дислокаций нет \\
\bottomrule
\end{longtable}

Примечания: HCP и OTHER приведены по матрице (98\,788 атомов); «beyond shell» — за пределами нормированной оболочки 1.3 от интерфейса. Термодинамика финала: eps0025 — $T = 286.2$\,K, $P = -2333$\,бар; eps0100 — $T = 286.8$\,K, $P = -2809$\,бар.

%======================= 8 =======================
\section{Анализ результатов}
\textbf{Подтверждённые дислокации.} В обоих вариантах DXA даёт ноль сегментов, нулевую длину и нулевую плотность.

\begin{quote}
\itshape DXA = 0 означает отсутствие подтверждённых дислокационных линий в данном масштабе расчёта.
\end{quote}

\textbf{HCP.} Доля HCP-атомов в матрице пренебрежимо мала (8 и 1 атом) и при усилении воздействия не растёт, а уменьшается. Развитой сетки дефектов упаковки нет.

\textbf{OTHER.} Высокая доля OTHER ($\approx 16$–$17\,\%$) объясняется структурным и тепловым фоном: атомы границы зерна, интерфейса Fe$_4$Al$_{13}$/Al, окрестности вакансий и атомы с большими тепловыми смещениями при 300\,K. Эта доля не растёт при четырёхкратном увеличении $\varepsilon_z$ ($16.34\,\%$ против $17.41\,\%$), то есть не является деформационной пластичностью.

\textbf{Почему OTHER нельзя называть дислокацией.} OTHER — это «всё нераспознанное», а не дислокация. Подтверждение дислокации даёт DXA по геометрии линии Бюргерса, а не счётчик нераспознанных атомов; DXA в обоих вариантах равен нулю. Итог: при масштабе 100k ни физический, ни усиленный вариант не дали подтверждённого зарождения дислокаций; наиболее вероятная причина — недостаточный размер модели для преодоления барьера зарождения.

%======================= 9 =======================
\section{План Stage C 1M}
Для проверки масштабного эффекта подготовлен и запущен масштабный расчёт.

\tabcap{Таблица 4. Параметры Stage C 1M.}
\begin{longtable}{p{0.46\textwidth} p{0.46\textwidth}}
\toprule
Параметр & Значение \\
\midrule
\endhead
Кейс & \code{C1\_1M\_nearGB\_vacancies\_medium\_eps0100} \\
Оценка числа атомов & 944\,812 \\
Воздействие & $\varepsilon_z = 0.0100$ (near-GB, вакансии) \\
Ожидаемая память GPU & 11.64 ГБ \\
Ожидаемое время счёта & $\approx 5$ суток \\
Оценка объёма данных & $\approx 3.98$ ГБ \\
Контрольная точка & 100\,000 шагов \\
Продолжение до 200k/250k & только вручную после анализа 100k \\
\bottomrule
\end{longtable}

На момент составления отчёта масштабный расчёт \textbf{запущен и находится в процессе выполнения} (run root \code{runs/stageC\_1M\_nearGB\_vacancies\_eps0100\_100k/20260616-173123}, PID 20944). Логика гейтов после 100k: при подтверждённой DXA — прицельный повтор окна события; при слабом предвестнике — продолжение до 200k/250k только после ревью; при чисто структурном фоне — остановка масштабирования и переход к ветке positive-control/seeded/cyclic.

%======================= 10 =======================
\section{Ограничения}
\begin{itemize}[leftmargin=1.6em,itemsep=1pt]
\item Масштаб 100k не гарантирует зарождение дислокаций: барьер зарождения и длина пробега дислокации могут не умещаться в ячейке $105\times105\times158$\,\AA.
\item Доля OTHER у границы/интерфейса не равна дислокации.
\item Масштаб $\approx$\,1 млн атомов необходим для проверки масштабного эффекта.
\item Per-atom напряжения — сравнительные вириальные оценки, а не абсолютные значения.
\item Обход соседних списков на GPU (\code{check no}) — валидированный локальный приём, а не исправление в исходниках LAMMPS.
\end{itemize}

%======================= 11 =======================
\section{Заключение}
\begin{enumerate}[leftmargin=1.8em,itemsep=1pt]
\item Построена и рассчитана реалистичная модель «Al-матрица + включение Fe$_4$Al$_{13}$» у границы зерна с вакансиями; магнитострикция реализована как собственная деформация включения.
\item В двух вариантах по 100\,000 шагов (физический $\varepsilon_z = 0.0025$ и усиленный $\varepsilon_z = 0.0100$) подтверждённые дислокации не зарегистрированы (DXA = 0).
\item Слабые структурные признаки (единичные HCP, высокий OTHER) объясняются фоном границы зерна, интерфейса и тепловыми смещениями и не зависят от деформации ожидаемым образом.
\item Обоснован и запущен масштабный расчёт $\approx$\,1 млн атомов (Stage C) для проверки масштабного эффекта; его результаты будут оформлены отдельным отчётом.
\end{enumerate}

%======================= 12 =======================
\section{Список источников}
\begin{enumerate}[leftmargin=1.8em,itemsep=1pt]
\item S. Plimpton. Fast Parallel Algorithms for Short-Range Molecular Dynamics // J. Comput. Phys. 117 (1995) 1–19.
\item M. I. Baskes. Modified embedded-atom potentials for cubic materials and impurities // Phys. Rev. B 46 (1992) 2727.
\item B. Jelinek et al. Modified embedded-atom method interatomic potentials for the Al–Si–Mg–Cu–Fe alloy system // Phys. Rev. B 85 (2012) 245102.
\item A. Stukowski. Visualization and analysis of atomistic simulation data with OVITO // Modelling Simul. Mater. Sci. Eng. 18 (2010) 015012.
\item A. Stukowski, K. Albe. Extraction of dislocations and non-dislocation crystal defects (DXA) // Modelling Simul. Mater. Sci. Eng. 18 (2010) 085001.
\end{enumerate}

%======================= 13 =======================
\section{Приложения}
\textbf{Приложение А. Малые подтверждающие артефакты (в репозитории):} \code{results/stageB\_focused\_100k\_summary.csv} и \code{.md}; \code{results/stageB\_focused\_100k\_artifacts/} (\code{analysis\_eps*.json}, \code{defect\_summary.csv}, \code{final\_report.md}, \code{eigenstrain\_build\_eps*.json}).

\textbf{Приложение Б. Stage C (launch/queue):} \code{results/stageC\_1M\_queue/} — preflight, оценка объёма, записи запуска, команды продолжения, effective config.

\textbf{Приложение В. Полные сырые данные:} только локально в \code{runs/} (см. \code{results/local\_raw\_data\_inventory.md}); в git не хранятся.

\end{document}
